\chapter{Macros and Threads}\label{ch:extensions}
\sloppy

The 1992 manual defines a language without macros and without threads.
The compiler to C adds both. They are \emph{extensions}: a program that uses them is
not a program of the 1992 language, and the interactive Common Lisp
implementation does not have them (it has its own extensions, described
in Chapter~\ref{ch:lisp}). Everything in this chapter applies to
programs compiled with \cmd{dylanc}.

\section{Macros}\label{sec:ext-macros}

\index{macros}
A \term{macro} is a function that runs while the program is being
compiled. It receives the forms of a call, unevaluated, and answers a new
form, its \term{expansion}, which the compiler compiles in place of the
call. Macros let you add new syntax to the language: a new looping form,
a definition form that writes several definitions, a form that evaluates
its argument in a special way.

The macros here are Common Lisp's kind. If you know \code{defmacro}, you
know them already: the parameters are bound to the argument forms, the
body builds the expansion, usually with backquote, and nothing protects
the expansion's names from the names around the call except the care of
the macro's writer and \code{gensym}.

The manual calls \code{and}, \code{or}, \code{cond}, \code{case},
\code{for} and the other syntax of Chapter~\ref{ch:control} macros. In
this system those are built into the compiler, not written with
\code{define-macro}, and the program cannot redefine them: the compiler
accepts a \code{define-macro} of one of their names without complaint,
but a call of the name still means the built-in form.

\subsection{\code{define-macro}}\label{sec:ext-define-macro}

\index{define-macro@\code{define-macro}}
\begin{quote}
\code{(define-macro \meta{name} \meta{parameter-list} \meta{form} ...)}
\end{quote}

\noindent at top level makes \meta{name}, a variable of the current
module, a macro. From then on, a call \code{(\meta{name}
\meta{argument} ...)} is replaced, while compiling, by the value of the
\meta{form}s, run with the parameters bound to the argument forms. The
parameter list is a method's: required parameters, \code{\#rest} and
\code{\#key} are all allowed.

\begin{dylan}
(define-macro my-unless (test #rest body)
  `(if ,test #f (begin ,@body)))

(my-unless (= 1 2)
  (print 'one)
  (print 'two))
(print (my-unless #t 'never))
\end{dylan}

\noindent prints

\begin{dylan}
one
two
#f
\end{dylan}

Here \code{test} is bound to the form \code{(= 1 2)} and \code{body} to
the list of the two \code{print} forms; the backquote template builds an
\code{if} form from them. Backquote follows the usual Lisp rules
(\code{,} inserts a value, \code{,@} splices a list, and backquotes may
nest); see Appendix~\ref{app:differences}.

Because the argument forms arrive unevaluated, a macro can use a form
both as data and as code:

\begin{dylan}
(define-macro show (form)
  `(format #t "~s => ~s~%" ',form ,form))

(show (+ 1 2))
(show (reverse '(a b c)))
\end{dylan}

\noindent prints

\begin{dylan}
(+ 1 2) => 3
(reverse (quote (a b c))) => (c b a)
\end{dylan}

A macro's body is ordinary Dylan, and it may compute its expansion any
way it likes; it need not use backquote, and its result need not even be
a list. Keyword parameters work as in a method:

\begin{dylan}
(define-macro scaled (#key (scale 1) value)
  (* scale value))
(print (list (scaled value: 5) (scaled scale: 3 value: 5)))
; prints (5 15)
\end{dylan}

Here the multiplication happens when the program is compiled; the
expansions are the constants \code{5} and \code{15}.

Macros may expand into uses of other macros, or of themselves; the
compiler expands again until no macro call is left. An expansion may also
be a definition, or a \code{begin} of several, at top level. A macro that
never stops expanding is reported (``macros expanded 500 deep --- one
never stops?'').

\subsection{Rules}

\begin{itemize}
\item \code{define-macro} must be at top level. Inside a method it is a
  compile error (``define-macro must be at top level'').
\item A macro must be defined before its first use in the program. A
  call that comes earlier is compiled as a call of a variable that is not
  defined: the compiler warns, and the call fails when it runs.
\item A name cannot be both a macro and a variable. \code{(define m 2)}
  after \code{(define-macro m () 1)}, or a use of \code{m} other than as
  the operator of a call, is a compile error (``m is a macro'').
\item A call of a macro must be a proper list, and must supply arguments
  its parameter list accepts (``macro m: wrong number of arguments'').
\item An error signalled by a macro's body is a compile error, reported
  at the call with the macro's name: \code{(define-macro m (a) (error "no
  good: \textasciitilde{}s" a))} and a call \code{(m 1)} report
  ``\code{macro m: no good: 1}''.
\end{itemize}

\subsection{Names in expansions: \code{gensym} and capture}
\label{sec:ext-hygiene}

\index{gensym@\code{gensym}}\index{macros!hygiene}
The names in an expansion mean what they mean \emph{where the macro is
used}. This is simple and predictable, but it means that the expansion
and the code around it can interfere. A macro that binds a variable of
its own risks capturing a variable of the same name at the call:

\begin{dylan}
(define-macro bad-swap! (a b)
  `(bind ((tmp ,a)) (set! ,a ,b) (set! ,b tmp)))

(bind ((tmp 1) (other 2))
  (bad-swap! tmp other)
  (print (list tmp other)))
; prints (1 2): nothing was swapped
\end{dylan}

The expansion binds a new \code{tmp}, and \code{(set! tmp other)} then
assigns that one, not the caller's. The cure is \code{(gensym)}, which
answers a fresh symbol that no program can write, for the expansion to
bind:

\begin{dylan}
(define-macro swap! (a b)
  (bind ((tmp (gensym "tmp")))
    `(bind ((,tmp ,a))
       (set! ,a ,b)
       (set! ,b ,tmp))))

(define tmp 1)
(define other 2)
(swap! tmp other)
(print (list tmp other))      ; prints (2 1)
\end{dylan}

\code{gensym} takes an optional string, the start of the new name, and
appends a number. The name only \emph{looks} like one a program could
write: \cmd{dylanc -E} shows it as \code{tmp1}, but it is not the symbol
\code{tmp1}, and nothing in the program can refer to it.

Capture can happen the other way too: an expansion that calls a function
by name gets whatever that name means at the call. If the caller has a
local variable of that name, the expansion uses it:

\begin{dylan}
(define-macro pair-of (x) `(list ,x ,x))
(print (pair-of 1))                     ; prints (1 1)
(bind ((list vector))
  (print (pair-of 2)))                  ; prints #(2 2)
\end{dylan}

\noindent There is no protection against this short of not binding, near
a macro's use, the names its expansions call. The built-in syntax ---
\code{for}, \code{case}, \code{with-lock} and the rest --- is not exposed
to either kind of capture: its expansions refer to the library's
variables directly and make their temporaries with names no program can
write.

A local variable or parameter with the name of a macro hides the macro,
as it would hide a module variable:

\begin{dylan}
(define-macro double (x) `(* 2 ,x))
(print (double 4))                              ; prints 8
(define apply-it (method (double) (double 4)))
(print (apply-it (method (n) (+ n 100))))       ; prints 104
\end{dylan}

\subsection{The compile-time environment and \code{define-for-macros}}
\label{sec:ext-compile-time}

\index{define-for-macros@\code{define-for-macros}}
Macro bodies run inside the compiler, in the interactive Common Lisp
implementation of Dylan (Chapter~\ref{ch:lisp}), which is loaded with
the compiler. They run in a \emph{compile-time environment} of their own,
separate from the program being compiled. It holds the whole library of
the interactive implementation, \code{gensym}, and whatever
\code{define-for-macros} has defined --- but nothing the program itself
defines. A macro body that calls one of the program's functions fails:

\begin{dylan}
(define helper (method (x) (list x x)))
(define-macro use-helper (x) (helper x))
(print (use-helper 3))
; error: macro use-helper: unbound variable helper
\end{dylan}

Functions and variables for macro bodies to use are defined with

\begin{quote}
\code{(define-for-macros \meta{form} ...)}
\end{quote}

\noindent at top level, which runs the \meta{form}s in the compile-time
environment while compiling, and leaves nothing in the compiled program.
A variable defined there keeps its value from one expansion to the next.
Here a helper builds the names of accessors for a macro that defines a
class:

\begin{dylan}
(define-for-macros
  (define accessor-name
    (method (class-name slot)
      (as <symbol>
          (concatenate (as <string> class-name) "-" (as <string> slot))))))

(define-macro define-record (name #rest slots)
  `(define-class ,name (<object>)
     ,@(map (method (s)
              `(,(accessor-name name s)
                init-keyword: ,(as <keyword> (as <string> s))))
            slots)))

(define-record point x y)
(bind ((p (make point x: 3 y: 4)))
  (print (list (point-x p) (point-y p))))   ; prints (3 4)
\end{dylan}

Two consequences of running macro bodies in the interactive
implementation are worth knowing. First, \code{gensym} and anything
\code{define-for-macros} defines exist only at compile time; a compiled
program cannot call \code{gensym}. Second, where the two implementations
differ, a macro body sees the interactive one: it has bignums, for
instance (Appendix~\ref{app:differences}).

\subsection{Macros and modules}\label{sec:ext-macro-modules}

A macro is a variable of the module it is defined in, and it is
exported, imported, renamed and excluded exactly as a variable is
(Chapter~\ref{ch:modules}). Here the module \code{macros} defines and
exports \code{repeat}, and \code{user} uses it:

\begin{dylan}
; macmods.dyl
(define-module macros (use dylan) (export repeat))
(define-module user (use dylan) (use macros))
\end{dylan}

\begin{dylan}
; macdef.dyl
module: macros

(define-macro repeat (n #rest body)
  (bind ((i (gensym)))
    `(dotimes (,i ,n) ,@body)))
\end{dylan}

\begin{dylan}
; macuse.dyl
module: user

(repeat 3 (print 'hi))
\end{dylan}

\begin{shell}
$ bin/dylanc -o macs macmods.dyl macdef.dyl macuse.dyl
$ ./macs
hi
hi
hi
\end{shell}

(The first lines of \file{macdef.dyl} and \file{macuse.dyl} are the
\code{module:} headers; the comments above them are not part of the
files.) The names in an expansion are looked up in the module where the
macro is \emph{used}, so \code{dotimes} must mean the same there as in
the macro's module --- here both use \code{dylan}.

\subsection{Seeing expansions: \cmd{dylanc -E}}\label{sec:ext-expand}

\index{dylanc@\cmd{dylanc}!-E@\cmd{-E}}
\cmd{dylanc -E} writes each use of the program's own macros, where it
occurs, and its expansion, and makes no C. It is the way to debug a
macro:

\begin{shell}
$ bin/dylanc -E m1.dyl
m1.dyl:4:1: (my-unless (= 1 2) (print 'one) (print 'two))
  => (if (= 1 2) #f (begin (print 'one) (print 'two)))
m1.dyl:7:8: (my-unless #t 'never)
  => (if #t #f (begin 'never))
m1.dyl:12:1: (show (+ 1 2))
  => (format #t "~s => ~s~%" '(+ 1 2) (+ 1 2))
m1.dyl:13:1: (show (reverse '(a b c)))
  => (format #t "~s => ~s~%" '(reverse '(a b c)) (reverse '(a b c)))
\end{shell}

Only the program's macros are shown, not the built-in syntax. Each step
of a macro that expands into another macro is shown separately.

\section{Threads}\label{sec:ext-threads}

\index{threads}
A compiled program can run several \term{threads} at once. Each Dylan
thread is a POSIX thread, so threads really do run in parallel on a
machine with several processors. The garbage collector is built for
threads and collects while they all allocate. Threads share the
program's module variables and every object they can reach; each has
its own stack, its own handlers and exits, and its own values of thread
variables (Section~\ref{sec:ext-thread-variables}).

The threads API is small. Table~\ref{tab:ext-threads} lists it; the
sections after it explain each part.

\begin{table}[htbp]
\centering
\small
\begin{tabular}{>{\raggedright\arraybackslash}p{0.47\textwidth}%
                >{\raggedright\arraybackslash}p{0.45\textwidth}}
\toprule
Form & What it does \\
\midrule
\code{(make-thread \meta{function})} & start a thread calling
  \meta{function}, which takes no arguments; answer the \code{<thread>} \\
\code{(join-thread \meta{thread})} & wait for \meta{thread} to end;
  answer its function's values \\
\code{(current-thread)} & the \code{<thread>} running \\
\code{(thread-yield)} & let other threads run; answer \code{\#f} \\
\code{(make <lock>)}, \code{(make <recursive-lock>)} & a new lock \\
\code{(acquire-lock \meta{lock})}, \code{(release-lock \meta{lock})} &
  hold and let go of a lock \\
\code{(with-lock (\meta{lock}) \meta{form} ...)} & the forms, holding
  the lock \\
\code{(make <condition-variable>)} & a new condition variable \\
\code{(wait-condition-variable \meta{cv} \meta{lock})} & release the
  lock, wait to be notified, acquire it again \\
\code{(notify-condition-variable \meta{cv})} & wake one waiting thread \\
\code{(notify-all-condition-variable \meta{cv})} & wake every waiting
  thread \\
\code{(make <atomic-integer> value: \meta{n})} & an integer cell for
  atomic updates (\meta{n} defaults to 0) \\
\code{(atomic-value \meta{a})} & its value \\
\code{(atomic-add! \meta{a} \meta{n})} & add \meta{n}; answer the new
  value \\
\code{(atomic-compare-and-swap! \meta{a} \meta{old} \meta{new})} & if
  it holds \meta{old}, store \meta{new} and answer \code{\#t}; else
  \code{\#f} \\
\code{(define-thread-variable \meta{name} \meta{value})} & a module
  variable with a value for each thread \\
\code{(dynamic-bind ((\meta{name} \meta{value}) ...) \meta{form} ...)} &
  the forms, with thread variables given new values in this thread \\
\code{<thread-error>}, \code{thread-error-thread},
  \code{thread-error-condition} & what \code{join-thread} signals when a
  thread ended with an unhandled condition \\
\bottomrule
\end{tabular}
\caption{The threads API.}\label{tab:ext-threads}
\end{table}

\subsection{Making and joining threads}\label{sec:ext-make-thread}

\index{make-thread@\code{make-thread}}\index{join-thread@\code{join-thread}}
\code{(make-thread \meta{function})} starts a new thread that calls
\meta{function} with no arguments, and answers at once with the
\code{<thread>} object. \code{(join-thread \meta{thread})} waits until
that thread's function returns, and answers the values it returned. A
thread may be joined more than once, and by any thread but itself; each
join answers the same values.

This program sums the integers below a million in four threads, each
summing a quarter, and adds the four results:

\begin{dylan}
(define sum-range
  (method (from below)
    (for ((i from (+ i 1))
          (s 0 (+ s i)))
         ((= i below) s))))

(define workers
  (map (method (k)
         (make-thread
           (method () (sum-range (* k 250000) (* (+ k 1) 250000)))))
       '(0 1 2 3)))

(print (reduce + 0 (map join-thread workers)))
; prints 499999500000
\end{dylan}

The threads run in whatever order the operating system chooses, but the
program's output does not depend on it: each thread computes alone, and
the main thread waits for all four before it adds. Write threaded
programs this way where you can --- threads that share nothing while they
work and are joined to collect their results.

A thread's function may return several values, and \code{join-thread}
answers them all:

\begin{dylan}
(define t (make-thread (method () (values 1 2 3))))
(print (bind ((a b c (join-thread t))) (list a b c)))   ; prints (1 2 3)
\end{dylan}

\code{(current-thread)} answers the thread that calls it. In the main
thread it answers a \code{<thread>} too, but the main thread cannot be
joined. \code{(thread-yield)} asks the operating system to run other
threads before this one continues; it answers \code{\#f}. A thread
prints as \code{\{thread\}}.

\subsection{When the program ends}\label{sec:ext-program-exit}

\index{threads!program exit}
The program ends when the main thread has run the last top-level form,
\emph{whatever its other threads are doing}. It does not wait for them;
they simply stop. This program prints \code{main-is-done} and exits with
status 0, although its other thread would loop forever:

\begin{dylan}
(define spinner
  (make-thread
    (method ()
      (bind-methods ((spin () (thread-yield) (spin)))
        (spin)))))
(print 'main-is-done)
\end{dylan}

Join every thread whose work must finish. Equally, an unhandled error in
the main thread ends the whole program at once, with the other threads.

\subsection{Unhandled conditions in threads}\label{sec:ext-thread-errors}

\index{thread-error@\code{<thread-error>}}
A serious condition that a thread does not handle ends \emph{that thread
only}. The program goes on. The condition is kept with the thread, and
\code{join-thread} on it signals a \code{<thread-error>} --- a subclass
of \code{<error>} --- in the joining thread. \code{thread-error-thread}
answers the thread that ended, and \code{thread-error-condition} the
condition that ended it:

\begin{dylan}
(define worker (make-thread (method () (error "disk ~a is full" 3))))

(handler-case (join-thread worker)
  ((<thread-error> condition: e)
   (bind ((original (thread-error-condition e)))
     (format #t "worker failed: ~a ~s~%"
             (condition-format-string original)
             (condition-format-arguments original))
     (print (id? (thread-error-thread e) worker)))))
\end{dylan}

\noindent prints

\begin{dylan}
worker failed: disk ~a is full (3)
#t
\end{dylan}

If the \code{<thread-error>} is not handled either, it ends the joining
thread in turn --- in the main thread, the program --- with a message
that carries the original one:

\begin{shell}
Dylan error: a thread ended with an unhandled condition: disk 3 is full
\end{shell}

A thread that is never joined ends silently when a condition is not
handled in it. Warnings are not serious conditions: an unhandled warning
in a thread is reported on the standard error stream, as in the main
thread, and the thread goes on.

\subsection{Locks}\label{sec:ext-locks}

\index{lock@\code{<lock>}}\index{recursive-lock@\code{<recursive-lock>}}
\index{with-lock@\code{with-lock}}
Threads that change shared data must take turns, or their changes will
be lost. A \term{lock} lets one thread at a time hold it. \code{(make
<lock>)} makes one; \code{(acquire-lock \meta{lock})} waits until no other
thread holds it and then holds it; \code{(release-lock \meta{lock})} lets
it go. Usually you will use

\begin{quote}
\code{(with-lock (\meta{lock}) \meta{form} ...)}
\end{quote}

\noindent which holds the lock while the forms run and releases it
however they are left --- by returning, by a non-local exit, or by an
error --- and answers the forms' values. Three threads each adding
100\,000 to a shared counter:

\begin{dylan}
(define counter 0)
(define counter-lock (make <lock>))
(define bump
  (method ()
    (dotimes (i 100000)
      (with-lock (counter-lock)
        (set! counter (+ counter 1))))))
(define threads (list (make-thread bump) (make-thread bump)
                      (make-thread bump)))
(for-each ((t threads)) () (join-thread t))
(print counter)                                  ; prints 300000
\end{dylan}

Without the lock, two threads can read the same old value of
\code{counter} and both store one more than it, and the total comes out
short.

A \code{<lock>} checks its use. Acquiring a \code{<lock>} the thread
already holds would wait forever, so it is an error instead
(``acquire-lock: the thread holds the lock already''), and so is
releasing a lock the thread does not hold (``release-lock: the thread
does not hold the lock''). A \code{<recursive-lock>} --- a subclass of
\code{<lock>} --- may be acquired again by the thread holding it, and is
released when it has been released as many times:

\begin{dylan}
(define r (make <recursive-lock>))
(print (with-lock (r) (with-lock (r) 'nested)))   ; prints nested
\end{dylan}

Neither kind of lock takes init-keywords.

\subsection{Condition variables}\label{sec:ext-condition-variables}

\index{condition-variable@\code{<condition-variable>}}
A thread that must wait until another thread has done something --- put
an item in a queue, say --- waits on a \term{condition variable}.
\code{(wait-condition-variable \meta{cv} \meta{lock})}, called while
holding \meta{lock}, which must be a \code{<lock>} (not a
\code{<recursive-lock>}), releases the lock and waits; when another
thread calls \code{(notify-condition-variable \meta{cv})}, one waiting
thread wakes, acquires the lock again, and returns.
\code{notify-all-condition-variable} wakes every waiting thread.

A wait can also end without a notification, and another thread may get
to the data first. So always wait in a loop that tests what is awaited,
with the lock held. Here a consumer thread takes ten numbers from a
queue that the main thread fills:

\begin{dylan}
(define lock (make <lock>))
(define not-empty (make <condition-variable>))
(define queue (make <deque>))

(define put!
  (method (item)
    (with-lock (lock)
      (push-last queue item)
      (notify-condition-variable not-empty))))

(define take!
  (method ()
    (with-lock (lock)
      (while (empty? queue)
        (wait-condition-variable not-empty lock))
      (pop queue))))

(define consumer
  (make-thread
    (method ()
      (for ((n 0 (+ n 1))
            (total 0 (+ total (take!))))
           ((= n 10) total)))))

(dotimes (i 10) (put! (* i i)))
(print (join-thread consumer))                  ; prints 285
\end{dylan}

However the two threads interleave, the consumer takes the ten squares
and answers their sum.

\subsection{Atomic integers}\label{sec:ext-atomic}

\index{atomic-integer@\code{<atomic-integer>}}
For a counter or a flag, a lock is more than is needed. An
\code{<atomic-integer>} holds an integer that threads can update without
one. \code{(make <atomic-integer> value: \meta{n})} makes one holding
\meta{n} (0 if \code{value:} is not given); \code{atomic-value} reads it;
\code{(atomic-add! \meta{a} \meta{n})} adds \meta{n} and answers the new
value, as one indivisible step; and \code{(atomic-compare-and-swap!
\meta{a} \meta{old} \meta{new})} stores \meta{new} only if \meta{a}
still holds \meta{old}, answering \code{\#t} if it did and \code{\#f} if
not.

\begin{dylan}
(define hits (make <atomic-integer>))
(define bump-atomically
  (method () (dotimes (i 100000) (atomic-add! hits 1))))
(do join-thread (list (make-thread bump-atomically)
                      (make-thread bump-atomically)))
(print (atomic-value hits))                          ; prints 200000
(print (atomic-compare-and-swap! hits 200000 0))     ; prints #t
(print (atomic-compare-and-swap! hits 200000 0))     ; prints #f
(print (atomic-add! hits 5))                         ; prints 5
\end{dylan}

The values must be integers, and an addition that overflows the 63-bit
range signals an error, as ordinary arithmetic does.

\subsection{Thread variables and \code{dynamic-bind}}
\label{sec:ext-thread-variables}

\index{define-thread-variable@\code{define-thread-variable}}
\index{dynamic-bind@\code{dynamic-bind}}
\begin{quote}
\code{(define-thread-variable \meta{name} \meta{value})}
\end{quote}

\noindent at top level defines a module variable that has a separate
value in each thread. Reading it, and assigning it with \code{set!},
affect only the current thread's value. The main thread starts with
\meta{value}, and so does every thread made after the definition has run
--- a new thread does \emph{not} inherit its creator's current value.

\begin{quote}
\code{(dynamic-bind ((\meta{name} \meta{value}) ...) \meta{form} ...)}
\end{quote}

\noindent gives each thread variable \meta{name} the new \meta{value} in
the current thread while the forms run, and restores the old value
however they are left. This is how such variables are normally used: as
settings that a computation and everything it calls can see, without
passing them as arguments, and without one thread's settings disturbing
another's.

\begin{dylan}
(define-thread-variable *indent* 0)

(define show
  (method (text)
    (format #t "~a~a~%" (make <string> size: *indent* fill: #\space) text)))

(define nested
  (method (depth)
    (show depth)
    (when (< depth 2)
      (dynamic-bind ((*indent* (+ *indent* 2)))
        (nested (+ depth 1))))))

(nested 0)
(print *indent*)
(print (join-thread (make-thread (method () (set! *indent* 99) *indent*))))
(print *indent*)
\end{dylan}

\noindent prints

\begin{dylan}
0
  1
    2
0
99
0
\end{dylan}

The other thread's \code{set!} changed only its own value.

\subsection{What threads may share}\label{sec:ext-thread-sharing}

Assigning a module variable or a slot is a single store of one machine
word, so a thread never sees a value half written. Beyond that, data
that threads share and change must be protected with locks. In
particular the library's collections --- tables, stretchy vectors,
deques and the rest --- do no locking of their own: two threads adding to
one table at once can leave it in a wrong state. The runtime's own shared state, such as the table of
symbols, is locked internally, and since the program's classes and
generic functions never change, method dispatch needs no locks at all.

\fussy
